\section{Legged and Wheeled-Legged Locomotion}
\label{sec:legged_wheeled_legged_locomotion}

Legged robots move through the environment directly through force contacts
exchange with the surrounding terrain. Unlike conventional wheeled
vehicles, they do not require a continuous and approximately flat path. Each
foot can be placed on a separate support surface, allowing the robot to cross
steps, gaps, debris, and other irregular terrain. This mobility makes legged
systems suitable for environments in which fixed infrastructure is unavailable
or ordinary vehicles cannot operate \cite{legged_optimization}.\\
\
This capability introduces a demanding control problem. A legged robot is a
floating-base mechanical system whose body is not directly actuated. Its motion
must be generated indirectly through the forces exchanged at the active
contacts. These contacts change during locomotion and are governed by
unilateral and friction constraints: a foot can push against the terrain but
cannot pull on it, while excessive tangential force may cause slipping. The
controller must coordinate joint motion, contact timing, and ground reaction
forces while maintaining balance and tracking the desired movement
\cite{legged_optimization}.\\
\
The complexity of locomotion also depends on the selected gait. Periodic gaits,
such as walking, trotting, and bounding, define recurring sequences of stance
and swing phases. During stance, the contacting limbs support and accelerate
the body. During swing, the remaining limbs are repositioned for the next
contact. Faster motions reduce the time available for feedback and may include
phases with a small support region or no ground contact. Aperiodic locomotion
provides more freedom to react to the terrain, but selecting contact timings
while planning the continuous motion makes the control problem considerably
harder \cite{legged_optimization,non_gaited}.\\
\
Model uncertainty further separates the predicted and actual motion. Impacts,
actuator limitations, delayed measurements, uncertain friction, and imperfect
terrain information cannot always be represented accurately. These effects are
particularly important when a controller relies on a reduced-order model, since
the neglected articulated dynamics still influence the physical robot. Robust
locomotion therefore requires continuous feedback and repeated adaptation of
the planned motion.\\
\
Wheeled-legged robots extend this mobility by placing wheels at the ends of
some or all limbs. On regular terrain, rolling avoids the repeated
acceleration and deceleration associated with conventional stepping and can
provide faster and more energy-efficient locomotion
\cite{wheeled_legged_survey,keep_rollin}. When the terrain becomes
discontinuous or unsuitable for rolling, the articulated legs can adjust the
body posture, negotiate obstacles, or initiate a stepping motion
\cite{rolling_deep}. The same platform can therefore use rolling for efficient
traversal and leg motion when the geometry of the environment requires it.\\
\
The mechanical arrangement of the wheels determines the available mobility.
Some platforms use steerable wheels, whereas others mount non-steerable wheels
at the ends of the legs. In the latter case, each contact can move along the
rolling direction but should not slide laterally. The controller must enforce
this non-holonomic condition while coordinating wheel velocity, leg
configuration, body motion, and contact forces
\cite{keep_rollin,wheeled_whole_body_mpc}. The wheels increase the set of
feasible motions, but also couple locomotion and posture control more tightly
than in a purely legged system.\\
\
Wheeled-legged locomotion can be divided into rolling, stepping, and hybrid
motions. During rolling, the wheels remain in contact with the ground and the
legs regulate the posture of the body. During stepping, one or more wheel
contacts are lifted and repositioned as ordinary feet. Hybrid locomotion
combines these modes within the same manoeuvre, allowing the robot to drive
across traversable regions and step over local obstacles
\cite{rolling_deep}. The contact sequence may be prescribed in advance or
adapted online according to the state of the robot and its kinematic limits
\cite{wheeled_whole_body_mpc}.\\
\
The addition of wheels does not remove the difficulties associated with legged
locomotion. Rolling contacts introduce further constraints, and transitions
between driving and stepping modify the support configuration. At higher
speeds, acceleration and cornering also redistribute the normal loads between
the contacts. This lateral load transfer reduces the available contact margin
and can increase the risk of wheel slip or rollover. On skid-steered platforms,
yaw is generated through differential wheel torques, further coupling
longitudinal motion, lateral dynamics, and body posture \cite{racing}.\\
\
Articulated legs can be used to mitigate these effects. By changing the
relative extension of the legs on either side of the body, a wheeled-legged
robot can regulate its roll angle and redistribute the contact loads. Physical
experiments on the Unitree Go2-W demonstrated this principle by using the legs
as an active suspension system and leaning the body into turns
\cite{racing}. This extends the role of the legs beyond obstacle traversal:
their configuration also affects stability and vehicle-like dynamics during
fast rolling.\\
\
Legged and wheeled-legged robots thus share the need to regulate an
underactuated floating body through constrained ground contacts. Wheels improve
continuous mobility on suitable terrain, but introduce rolling constraints,
additional actuation, and high-speed effects that are absent or less prominent
in conventional stepping. These properties motivate controllers that can plan
body motion and contact forces online while remaining sufficiently efficient
for closed-loop operation.
